
\documentclass{article}

\usepackage[usenames]{color}
\usepackage{colortbl}
\usepackage[unicode, pdftex]{hyperref}
\usepackage[T2A]{fontenc}
\usepackage[utf8]{inputenc}
\usepackage[english,russian]{babel}
\usepackage{amsmath,amsthm,amssymb,amsfonts,mathtools}
\usepackage{mathtext}

\begin{document}
\hypertarget{n0}{}
\begin{center}
    \textbf{\Large Ковальский И. В. (Б22-514)}
    \end{center}


\begin{center}
\textbf{\Large Вариант 10}
\end{center}


\begin{center}
    \textbf{\Large Оглавление(быстрый поиск)}
    \end{center}
 \hyperlink{n1}{\textcolor{blue}{\fbox{Номер 1}}}
\hyperlink{n2}{\textcolor{blue}{\fbox{Номер 2}}}
\hyperlink{n3}{\textcolor{blue}{\fbox{Номер 3}}}
\hyperlink{n4}{\textcolor{blue}{\fbox{Номер 4}}}
\hyperlink{n5}{\textcolor{blue}{\fbox{Номер 5}}}
\hyperlink{n6}{\textcolor{blue}{\fbox{Номер 6}}}
\hyperlink{n7}{\textcolor{blue}{\fbox{Номер 7}}}
\hyperlink{n8}{\textcolor{blue}{\fbox{Номер 8}}}
\hypertarget{n1}{}
\section{}

\begin{center}
    \textbf{\Large 1.10}
\end{center}


\textbf{\Large Задание}


Свести двойной интеграл 
$\iint  f(x,y) \,dx \,dy$
к повторному, расставить пределы интегрирования в полярных координатах в том и другом порядке, если область G задана следующим образом:
\begin{center}
G: $x^2 + (y-1)^2 <= 1$  , x >= 0
\end{center}

\textbf{\Large Решение}
\begin{center}
Функция G - внутренняя часть окр. с центром в точке (0,1) и R = 1
\end{center}

\begin{equation*}
\begin{cases}
    x = r 
\cos \varphi
\\
y = r
\sin \varphi
\end{cases}
\end{equation*}

Подставляя новые координаты в функцию G получаем
r = 2
$\sin \varphi$

Из условия задачи, что x >= 0 и r >= 0 находим отрезок, в котором лежит
$\varphi$:

\begin{center}
    $\varphi$ =
$[0,\frac{\pi}{2}]$
\end{center}

Радиус r лежит в следующих пределах:

\begin{center}
    r =
$[0,2\sin \varphi]$
\end{center}

I = $\int_{0}^{\frac{\pi}{2}}  \,d\varphi \int_{0}^{2\sin \varphi} rf(r\cos \varphi,r\sin \varphi) \,dr$

Меняя предел интегрирования получаем:
I = $\int_{0}^{2} r  \,dr \int_{0}^{\arcsin \frac{r}{2}} f(r\cos \varphi,r\sin \varphi) \,d\varphi$

\textbf{\Large Ответ}

I = $\int_{0}^{\frac{\pi}{2}}  \,d\varphi \int_{0}^{2\sin \varphi} rf(r\cos \varphi,r\sin \varphi) \,dr = \int_{0}^{2} r  \,dr \int_{0}^{\arcsin \frac{r}{2}} f(r\cos \varphi,r\sin \varphi) \,d\varphi$
\hypertarget{n2}{}
\[\hyperlink{n0}{\textcolor{blue}{\fbox{Вернуться в начало}}}\]
\section{}

\begin{center}
    \textbf{\Large 2.10}
\end{center}

\textbf{\Large Задание}

Пользуясь цилиндрическими или сферическими координатами, найти объем тела, ограниченного линиями:
\begin{center}
\[\frac{x^2}{2^2}+\frac{y^2}{3^2} = 1, z = xy, z = 0, x > 0, y > 0\]
\end{center}

\textbf{\Large Решение}

Перейдем к цилиндрическим координатам:
\begin{equation*}
    \begin{cases}
        x = r 
    \cos \varphi
    \\
    y = r
    \sin \varphi
    \\
    z = h
    \\
    J = r
    \end{cases}
    \end{equation*}

Подставив их в первое уравнение получаем выражение для r:
\begin{center}
    \[r = \frac{6}{\sqrt{5\cos ^ 2 \varphi + 4} }\]
 \end{center}

Из условия задачи, что x > 0 и y > 0 находим отрезок, в котором лежит
$\varphi$:

\begin{center}
    $\varphi$ =
$(0,\frac{\pi}{2})$
\end{center}

Из уравнения z = xy находим h:
\begin{center}
    $h=r^2\cos \varphi\sin \varphi$
    \end{center}

Считаем следующий интеграл:
\[I = \int_{0}^{\frac{\pi}{2}}  \,d\varphi \int_{0}^{\frac{6}{\sqrt{5\cos ^ 2 \varphi + 4} }} r  \,dr \int_{0}^{r^2\cos \varphi\sin \varphi}  \,dh\]

\[I = \int_{0}^{\frac{\pi}{2}}  \,d\varphi \int_{0}^{\frac{6}{\sqrt{5\cos ^ 2 \varphi + 4} }} r  \,dr \int_{0}^{r^2\cos \varphi\sin \varphi}  \,dh = \int_{0}^{\frac{\pi}{2}}  \,d\varphi \int_{0}^{\frac{6}{\sqrt{5\cos ^ 2 \varphi + 4} }} r ^ 3 cos \varphi\sin \varphi \,dr =\] 

= 162 $\int_{0}^{\frac{\pi}{2}}   \frac{4\sin 2\varphi}{(5\cos 2\varphi + 13)^ 2} \,d\varphi = \frac{9}{2}$

\textbf{\Large Ответ}

$\frac{9}{2}$
\hypertarget{n3}{}
\[\hyperlink{n0}{\textcolor{blue}{\fbox{Вернуться в начало}}}\]
\section{}

\begin{center}
    \textbf{\Large 3.10}
\end{center}

\textbf{\Large Задание}

Вычислить криволинейный интеграл первого рода скалярного поля f вдоль кривой C:
\begin{center}
$\int\limits_{C} (x - y) \,dS$   
\end{center}     

\begin{center}
$C: x ^ 2 + y^2 = 2y$
\end{center}     

\textbf{\Large Решение}

Контур C представляет собой окружность

Сделаем замену: 
\begin{equation*}
    \begin{cases}
        x = 2
    \cos \varphi \sin \varphi
    \\
    y = 2
    \cos ^ 2 \varphi
   
    \end{cases}
    \end{equation*}



        \[dS = \sqrt{\varphi\backprime ^ (2) (t) + \psi\backprime ^ (2) (t)  }\,dt  = 2 \,dt\]

        \[I = \int_{0}^{2\pi} (2\cos \varphi \sin \varphi - 2\cos ^ 2 \varphi)2 \,dt = 2 \int_{0}^{2\pi} (\sin 2\varphi - \cos 2\varphi - 1) \,dt = \]

        \[= 2(\frac{1}{2} - 2\pi - \frac{1}{2}) = 4\pi\]

\textbf{\Large Ответ}

$4\pi$
\hypertarget{n4}{}
\[\hyperlink{n0}{\textcolor{blue}{\fbox{Вернуться в начало}}}\]
\section{}

\begin{center}
    \textbf{\Large 4.10}
\end{center}

\textbf{\Large Задание}

Вычислить поверхностный интеграл первого рода скалярного поля f по поверхности S:
\begin{center}
$\iint\limits_{S} \left\lvert xy \right\rvert z \,dS$
\end{center}

\begin{center}
    $S: z = x^2 + y^2, z <= 1$
\end{center} 

\textbf{\Large Решение}

\[\frac{\partial  z}{\partial  x} = 2x, \frac{\partial  z}{\partial  y} = 2y\]

\[\,dS = \sqrt{1 + (\frac{\partial  z}{\partial  x})^2 + (\frac{\partial  z}{\partial  y})^2} \,dx \,dy = \sqrt{1 + 4x^2 + 4y^2}\,dx \,dy\]

\[\iint\limits_{S} \left\lvert xy \right\rvert z \,dS = \iint \left\lvert xy \right\rvert (x^2 + y^2) \sqrt{1 + 4x^2 + 4y^2} \,dx \,dy\]

Перейдем к полярным координатам:

\begin{equation*}
    \begin{cases}
        x = r 
    \cos \varphi
    \\
    y = r
    \sin \varphi
    \end{cases}
    \end{equation*}

Из условий $z = x^2 + y^2, z <= 1$ находим, что r <= 1 

Получаем следующие пределы интегрирования:

\begin{center}
    $\varphi$ =
$[0,2\pi]$
\end{center}

\begin{center}
    r =
$[0,1]$
\end{center}

\[\iint \left\lvert xy \right\rvert (x^2 + y^2) \sqrt{1 + 4x^2 + 4y^2} \,dx \,dy  = \iint \left\lvert r^2 \cos \varphi \sin \varphi \right\rvert r^3 \sqrt{1+4r^2} \,dr \,d\varphi =\] 

= $\frac{125\sqrt{5} - 1}{420}$

\textbf{\Large Ответ}

$\frac{125\sqrt{5} - 1}{420}$
\hypertarget{n5}{}
\[\hyperlink{n0}{\textcolor{blue}{\fbox{Вернуться в начало}}}\]
\section{}

\begin{center}
    \textbf{\Large 5.10}
\end{center}

\textbf{\Large Задание}

Проверить, является ли веткорное поле $\overrightarrow{F}$ потенциальным; если является, то найти потенциал и работу поля при перемещении из точки (1;2;1) в точку (4;3;2):
\begin{center}
$\overrightarrow{F} = (y^4 - 6xz^3)\overrightarrow{i} + (3y^2z^2 + 4xy^3)\overrightarrow{j} + (2y^3z-9x^2z^2)\overrightarrow{k}$
\end{center}

\textbf{\Large Решение}

Найдем ротор $rot \overrightarrow{F}$:

\[
    rot \overrightarrow{F} =
    \begin{vmatrix}
        \overrightarrow{i} & \overrightarrow{j} & \overrightarrow{k}\\
        \frac{\partial }{\partial  x} & \frac{\partial }{\partial  y} & \frac{\partial }{\partial  z}\\
      F_{x} & F_{y}& F_{z}
    \end{vmatrix}
 = (6y^2z - 6y^2z)\overrightarrow{i} - (-18xz^2 + 18xz^2) \overrightarrow{j} + (4y^3 - 4y^3) \overrightarrow{k} = 0\]
 
 Значит поле $\overrightarrow{F}$ - потенциально

 \[
\frac{\partial  U}{\partial  x} = y^4 - 6xz^3,
 \frac{\partial  U}{\partial  y} = 3y^2z^2 + 4xy^3,
  \frac{\partial  U}{\partial  z} = 2y^3z - 9x^2z^2
\]

\[U(x,y,z) = \int (y^4 - 6xz^3) \,dx + C_{1}(y,z) = xy^4 - 3x^2z^3 + \varphi(y,z)\]

\[\frac{\partial  U}{\partial  y} = 4xy^3 + \frac{\varphi(y,z)}{\partial  y} = 4xy^3 + 3z^2y^2 \implies \frac{\partial  \varphi}{\partial  y} = 3y^2z^2\]

\[\varphi(y,z) = \int 3z^2y^2 \,dy + C_{2}(z) = y^3z^2 + \psi(z)\]

\[\frac{\partial  U}{\partial  z} = \frac{\partial }{\partial  z}(xy^4 - 3x^2z^3 + y^3z^2 + \psi(z)) \implies \]

\[\implies -9x^2z^2 + 2y^3z + \psi\backprime (z) = 2y^3z - 9x^2z^2 \implies \psi\backprime (z) = 0\]
Следовательно, $\psi(z) = C = const$
\[U(x,y,z) = xy^4 - 3x^2z^3 + y^3z^2 + C\]
Работа поля при перемещении из точки (1;2;1) в точку (4;3;2):
\[A = U(4,3,2) - U(1,2,1) = 27\]
\textbf{\Large Ответ}

\[U(x,y,z) = xy^4 - 3x^2z^3 + y^3z^2 + C\] \[A = 27\]
\hypertarget{n6}{}
\[\hyperlink{n0}{\textcolor{blue}{\fbox{Вернуться в начало}}}\]
\section{}

\begin{center}
    \textbf{\Large 6.10}
\end{center}

\textbf{\Large Задание}

Найти поток векторного поля $\overrightarrow{F}$ через внешнюю сторону поверхности S:
\[\overrightarrow{F} = x^2\overrightarrow{i} + y^2\overrightarrow{j} + z^2\overrightarrow{k}\]
S - внешняя сторона сферы
\[(x-1)^2 + (y-1)^2 + z^2 = R^2 \]
\textbf{\Large Решение}
\[div \overrightarrow{F} = 2x + 2y + 2z\]
По формуле Остроградского-Гаусса
\[\Pi = \iint\limits_{\delta} \overrightarrow{F} \,d\overrightarrow{\delta} = \iiint\limits_{V} div \overrightarrow{F} \,dV = \iiint\limits_{V} (2x + 2y + 2z) \,dV = 2\iiint\limits_{V} (x + y + z) \,dV\]
Перейдем к сферическим координатам:
\begin{equation*}
    \begin{cases}
        x = 1 + r 
    \cos \varphi \sin \theta
    \\
    y = 1 + r
    \sin \varphi \sin \theta
    \\
    z = r \cos\theta
    \\
    J = r^2\sin\theta
    \end{cases}
    \end{equation*}

\begin{center}
    $\varphi$ =
    $[0,2\pi]$
\end{center}

\begin{center}
        $\theta$ =
    $[0,\pi]$
\end{center}

\begin{center}
        r =
    $[0,R]$
\end{center}

Произведя замену считаем следующий интеграл:
\[2\int_{0}^{2\pi}  \,d\varphi \int_{0}^{\pi}  \,d\theta \int_{0}^{R} (1 + r\cos\varphi \sin\theta + 1 + r\sin\varphi\sin\theta + r\cos\theta)r^2\sin\theta  \,dr = \frac{16\pi R^3}{3}\]
\textbf{\Large Ответ}

$\frac{16\pi R^3}{3}$




\hypertarget{n7}{}
\[\hyperlink{n0}{\textcolor{blue}{\fbox{Вернуться в начало}}}\]
\section{}

\begin{center}
    \textbf{\Large 7.10}
\end{center}

\textbf{\Large Задание}

Используя формулу Гаусса-Остроградского, найти поток векторного поля $\overrightarrow{F}$ через всю границу области S.
\[\overrightarrow{F} = 4x\overrightarrow{i} - 2y\overrightarrow{j} - z\overrightarrow{k}\]
\[S:3x + 2y = 12, 3x + y = 6, y = 0, x + y + z = 6, z = 0 \]
\textbf{\Large Решение}

$div \overrightarrow{F} = 4 - 2 - 1 = 1$

\[\iint\limits_{S}(\overrightarrow{F} * \overrightarrow{n})  \,dS  = \iiint\limits_{F}1 \,dx\,dy\,dz = \int_{0}^{6}  \,dy \int_{\frac{6 - y}{3}}^{\frac{12 - 2y}{3}}  \,dx \int_{0}^{6 - x - y}  \,dz  \]
Считая данный интеграл получаем:
\[\frac{1}{6}\int_{0}^{6} (y - 6) ^ 2 \,dy = -12\]
\textbf{\Large Ответ:}
-12
\hypertarget{n8}{}
\[\hyperlink{n0}{\textcolor{blue}{\fbox{Вернуться в начало}}}\]
\section{}

\begin{center}
    \textbf{\Large 8.10}
\end{center}

\textbf{\Large Задание}

Используя формулу Стокса, вычислить циркуляцию векторного поля $\overrightarrow{F}$ вдоль кривой $\gamma$

\[\overrightarrow{F} = y\overrightarrow{i} - x\overrightarrow{j} + z^2\overrightarrow{k}\]
\[\gamma: x^2 + y^2 = 1, z = 4\]
\textbf{\Large Решение}

z = 4 - поверхность S, через которую будем искать поток ротора.

Единичный вектор - (0,0,1)

\[
    rot \overrightarrow{F} =
    \begin{vmatrix}
        \overrightarrow{i} & \overrightarrow{j} & \overrightarrow{k}\\
        \frac{\partial }{\partial  x} & \frac{\partial }{\partial  y} & \frac{\partial }{\partial  z}\\
      P & Q & R
    \end{vmatrix}
 = (0)\overrightarrow{i} - (0) \overrightarrow{j} + (-1 - 1) \overrightarrow{k} = -2\overrightarrow{k}\]

 \[\oint\limits_{\gamma} \overrightarrow{F} \cdot \,d\overrightarrow{r} = -2 \iint\limits_{S} \,dS         \]
Т.к. z = 4, то $ \,dS = \sqrt{1 + (\frac{\partial  z}{\partial  x})^2 + (\frac{\partial  z}{\partial  y})^2} \,dx \,dy = \sqrt{1} \,dx \,dy$
\[-2 \iint\limits_{S} \,dS = -2 \iint\limits_{D} \,dx \,dy = -2\pi r^2 = -2\pi\]

Поскольку область D: $x^2 + y^2 = 1$ представляет собой окружность радиуса 1 с центром в т. (0,0)

\textbf{\Large Ответ:}
-$2\pi$

\[\hyperlink{n0}{\textcolor{blue}{\fbox{Вернуться в начало}}}\]





























































































































\end{document}
